\chapter{Lecture 14}

\section{Brauer--Nesbitt and Decomposition Numbers}

Recall a \(p\)-modular system \((F,R,k)\), where \(R\) is a complete discrete valuation ring
with uniformizer \(\pi\), fraction field \(F\), and residue field \(k=R/(\pi)\).

\begin{theorem}[Brauer--Nesbitt]
Let \((F,R,k)\) be a \(p\)-modular system, and let \(G\) be a finite group.
Let \(U\) be a finitely generated \(FG\)-module.
Let \(L_1,L_2\) be full \(RG\)-lattices in \(U\).
Then
\[
L_1/\pi L_1
\qquad\text{and}\qquad
L_2/\pi L_2
\]
have the same composition factors, with the same multiplicities, as \(kG\)-modules.
\end{theorem}

\begin{proof}
First observe that \(L_1+L_2\) is also a full \(RG\)-lattice in \(U\).
Thus, replacing \(L_1,L_2\) if necessary, we may assume
\[
L_1\subseteq L_2.
\]

As \(R\)-modules, \(L_1\) and \(L_2\) are free of the same rank. Hence \(L_2/L_1\)
is a torsion \(R\)-module. Thus \(L_2/L_1\) has a finite composition series as an
\(R\)-module, and hence we may reduce to the case where \(L_1\) is a maximal
\(RG\)-submodule of \(L_2\).

Then \(L_2/L_1\) is simple. Since \(L_2/L_1\) is torsion over the DVR \(R\), we have
\[
\pi L_2\subseteq L_1.
\]
Thus
\[
\pi L_1\subseteq \pi L_2\subseteq L_1\subseteq L_2.
\]

It remains to show that
\[
L_2/L_1 \cong \pi L_2/\pi L_1.
\]
Indeed, define
\[
\varphi:L_2\longrightarrow \pi L_2/\pi L_1,
\qquad
x\longmapsto \pi x+\pi L_1.
\]
This is a surjective \(RG\)-homomorphism, and
\[
\ker(\varphi)=L_1.
\]
Therefore
\[
L_2/L_1\cong \pi L_2/\pi L_1.
\]

Using the chains
\[
\pi L_1\subseteq \pi L_2\subseteq L_1
\]
and
\[
\pi L_2\subseteq L_1\subseteq L_2,
\]
the above isomorphism shows that \(L_1/\pi L_1\) and \(L_2/\pi L_2\) have the same
composition factors with the same multiplicities.
\end{proof}

\begin{definition}[Decomposition matrix]
Assume \((F,R,k)\) is splitting for \(G\).
The decomposition matrix \(D=(d_{TS})\) is the matrix whose rows are indexed by
the simple \(FG\)-modules \(T\), and whose columns are indexed by the simple
\(kG\)-modules \(S\). Its entries are
\[
d_{TS}
=
\left[T_0/\pi T_0:S\right],
\]
where \(T_0\) is a full \(RG\)-lattice in \(T\), and \(S\) is a simple \(kG\)-module.

The integer \(d_{TS}\) is called a decomposition number.
\end{definition}

\section{Grothendieck Groups and the Cartan Triangle}

Assume \((F,R,k)\) is splitting for \(G\), and assume \(R\) is complete.

Define \(G_0(FG)\) to be the free abelian group with basis the isomorphism classes of
simple \(FG\)-modules. Similarly, define \(G_0(kG)\) to be the free abelian group with
basis the isomorphism classes of simple \(kG\)-modules.

Define \(K_0(kG)\) to be the free abelian group with basis the isomorphism classes of
indecomposable projective \(kG\)-modules.

If \(T\) is a simple \(FG\)-module, write \([T]\) for the corresponding basis element of
\(G_0(FG)\). If \(S\) is a simple \(kG\)-module, write \([S]\) for the corresponding basis
element of \(G_0(kG)\). If \(P\) is an indecomposable projective \(kG\)-module, write
\([P]\) for the corresponding basis element of \(K_0(kG)\).

If \(U\) is a \(kG\)-module with composition factors
\[
S_1,\dots,S_r
\]
with multiplicities
\[
n_1,\dots,n_r,
\]
we write
\[
[U]=n_1[S_1]+\cdots+n_r[S_r]\in G_0(kG).
\]

Similarly, for an \(FG\)-module \(V\), if
\[
V\cong T_1^{n_1}\oplus\cdots\oplus T_s^{n_s},
\]
where the \(T_i\) are simple \(FG\)-modules, we write
\[
[V]=n_1[T_1]+\cdots+n_s[T_s]\in G_0(FG).
\]

If \(P\) is a projective \(kG\)-module and
\[
P\cong P_1^{n_1}\oplus\cdots\oplus P_s^{n_s},
\]
where the \(P_i\) are indecomposable projective \(kG\)-modules, then
\[
[P]=n_1[P_1]+\cdots+n_s[P_s]\in K_0(kG).
\]

Consider the following Cartan triangle:
\[
\begin{tikzcd}
& G_0(FG) \arrow[dr,"d"] & \\
K_0(kG) \arrow[ur,"e"] \arrow[rr,"c"'] && G_0(kG).
\end{tikzcd}
\]

If \(P_S\) is an indecomposable projective \(kG\)-module, then
\[
e([P_S])=\left[F\otimes_R \widehat P_S\right]\in G_0(FG),
\]
where \(\widehat P_S\) is an \(RG\)-lattice lifting \(P_S\). Recall that \(\widehat P_S\)
is unique up to isomorphism.

If \(V\) is a simple \(FG\)-module containing a full \(RG\)-lattice \(V_0\), define
\[
d([V])=[V_0/\pi V_0]\in G_0(kG).
\]
With respect to the bases above, the matrix of \(d\) is the transpose of the decomposition
matrix \(D\).

Let \(P_S\) be the indecomposable projective \(kG\)-module which is the projective cover
of a simple \(kG\)-module \(S\). Define
\[
c([P_S])=[P_S].
\]
Thus \(c:K_0(kG)\to G_0(kG)\) sends a projective module to its class as a \(kG\)-module.
If
\[
[P_T]=\sum_{S} C_{ST}[S],
\]
where \(S,T\) run through simple \(kG\)-modules, then the matrix of \(c\) is the Cartan
matrix
\[
C=(C_{ST}).
\]

\begin{theorem}
With the notation above,
\[
c=d\circ e.
\]
\end{theorem}

\begin{proof}
Let \(P_S\) be an indecomposable projective \(kG\)-module. Then
\[
c([P_S])=[P_S],
\]
while
\[
e([P_S])=\left[F\otimes_R \widehat P_S\right].
\]
Choose the full \(RG\)-lattice \(\widehat P_S\) in \(F\otimes_R\widehat P_S\).
Reducing modulo \(\pi\), we get
\[
d\left(\left[F\otimes_R \widehat P_S\right]\right)=[P_S].
\]
Hence
\[
(d\circ e)([P_S])=[P_S]=c([P_S]).
\]
Since the classes \([P_S]\) form a basis of \(K_0(kG)\), it follows that
\[
c=d\circ e.
\]
\end{proof}

\section{Hom-Lattices}

\begin{proposition}
Let \((F,R,k)\) be a \(p\)-modular system. Let \(U,V\) be \(FG\)-modules containing full
\(RG\)-lattices \(U_0,V_0\), respectively.

\begin{enumerate}
\item The \(R\)-module
\[
\Hom_{RG}(U_0,V_0)
\]
is a full \(RG\)-lattice in
\[
\Hom_{FG}(U,V).
\]

\item If \(U_0\) is a projective \(RG\)-lattice, then
\[
\Hom_{RG}(U_0,V_0)/\pi \Hom_{RG}(U_0,V_0)
\cong
\Hom_{RG}(U_0,V_0/\pi V_0)
\]
and
\[
\Hom_{RG}(U_0,V_0/\pi V_0)
\cong
\Hom_{kG}(U_0/\pi U_0,V_0/\pi V_0).
\]
\end{enumerate}
\end{proposition}

\begin{proof}
First, \(\Hom_{RG}(U_0,V_0)\) may be viewed as a subset of \(\Hom_{FG}(U,V)\).
Choose \(R\)-bases
\[
u_1,\dots,u_r
\qquad\text{and}\qquad
v_1,\dots,v_s
\]
of \(U_0\) and \(V_0\), respectively. These are also \(F\)-bases of \(U\) and \(V\).
Any \(RG\)-homomorphism \(U_0\to V_0\) is represented by a matrix over \(R\), and
therefore determines an \(FG\)-homomorphism \(U\to V\).

We have
\[
\Hom_{RG}(U_0,V_0)\subseteq \Hom_R(U_0,V_0)\cong R^{rs}.
\]
Since \(R\) is a PID, this submodule is \(R\)-free.

It remains to show fullness. Let
\[
\phi:U\to V
\]
be an \(FG\)-homomorphism. Write
\[
\phi(u_i)=\sum_j \lambda_{ji}v_j,
\qquad
\lambda_{ji}\in F.
\]
Choose \(a\in R\setminus\{0\}\) such that
\[
a\lambda_{ji}\in R
\]
for all \(i,j\). Then
\[
a\phi:U_0\to V_0
\]
is an \(RG\)-homomorphism. Therefore
\[
\Hom_{RG}(U_0,V_0)
\]
spans \(\Hom_{FG}(U,V)\) over \(F\), and hence is a full lattice.

Now assume \(U_0\) is projective. Then
\[
\Hom_{RG}(U_0,\pi V_0)=\pi\Hom_{RG}(U_0,V_0).
\]
Consider the natural map
\[
\Hom_{RG}(U_0,V_0)
\longrightarrow
\Hom_{RG}(U_0,V_0/\pi V_0).
\]
Since \(U_0\) is projective, this map is surjective. Its kernel is
\[
\Hom_{RG}(U_0,\pi V_0)=\pi\Hom_{RG}(U_0,V_0).
\]
Hence
\[
\Hom_{RG}(U_0,V_0)/\pi\Hom_{RG}(U_0,V_0)
\cong
\Hom_{RG}(U_0,V_0/\pi V_0).
\]

For the second isomorphism, every morphism
\[
U_0\to V_0/\pi V_0
\]
contains \(\pi U_0\) in its kernel, and therefore factors uniquely through
\[
U_0/\pi U_0.
\]
Thus
\[
\Hom_{RG}(U_0,V_0/\pi V_0)
\cong
\Hom_{kG}(U_0/\pi U_0,V_0/\pi V_0).
\]
\end{proof}

\begin{corollary}
Suppose \(U_0,V_0\) are full \(RG\)-lattices in \(FG\)-modules \(U,V\), respectively.
Assume \(U_0\) is projective. Then
\[
\operatorname{dim}_F\Hom_{FG}(U,V)
=
\operatorname{dim}_k\Hom_{kG}(U_0/\pi U_0,V_0/\pi V_0).
\]
\end{corollary}

\begin{proof}
Both sides are equal to the \(R\)-rank of \(\Hom_{RG}(U_0,V_0)\).
\end{proof}

\section{Brauer Reciprocity and the Cartan Matrix}

\begin{theorem}[Brauer reciprocity]
Assume \((F,R,k)\) is splitting for \(G\), and \(R\) is complete.
Let \(S\) be a simple \(kG\)-module, and let \(T\) be a simple \(FG\)-module containing
a full \(RG\)-lattice \(T_0\).

Let \(P_S\) be the projective cover of \(S\), and let \(\widehat P_S\) be its lift to an
\(RG\)-lattice. Then the multiplicity of \(T\) in
\[
F\otimes_R\widehat P_S
\]
is equal to the multiplicity of \(S\) as a composition factor of
\[
T_0/\pi T_0.
\]
Equivalently,
\[
\left[F\otimes_R\widehat P_S:T\right]
=
\left[T_0/\pi T_0:S\right].
\]
\end{theorem}

\begin{proof}
By the preceding corollary,
\[
\operatorname{dim}_F\Hom_{FG}\left(F\otimes_R\widehat P_S,T\right)
=
\operatorname{dim}_k\Hom_{kG}\left(P_S,T_0/\pi T_0\right).
\]

The left-hand side equals
\[
\operatorname{dim}_F\operatorname{End}_{FG}(T)\cdot
\left[F\otimes_R\widehat P_S:T\right].
\]
The right-hand side equals
\[
\operatorname{dim}_k\operatorname{End}_{kG}(S)\cdot
\left[T_0/\pi T_0:S\right].
\]
Since \((F,R,k)\) is splitting, we have
\[
\operatorname{dim}_F\operatorname{End}_{FG}(T)=1,
\qquad
\operatorname{dim}_k\operatorname{End}_{kG}(S)=1.
\]
Therefore
\[
\left[F\otimes_R\widehat P_S:T\right]
=
\left[T_0/\pi T_0:S\right].
\]
\end{proof}

With respect to the standard bases, the matrix of \(e\) is \(D\), and the matrix of \(d\)
is \(D^T\), where \(D\) is the decomposition matrix.

\begin{corollary}
Assume \((F,R,k)\) is splitting for \(G\), and \(R\) is complete. Then the Cartan matrix
satisfies
\[
C=D^T D.
\]
In particular, \(C\) is symmetric.
\end{corollary}

\begin{proof}
Since
\[
c=d\circ e,
\]
and the matrices of \(d\) and \(e\) are \(D^T\) and \(D\), respectively, the matrix of \(c\)
is
\[
C=D^T D.
\]
\end{proof}

\section{Integral Elements}

Let \(S\) be a commutative ring with \(1\), and let \(R\subseteq S\) be a subring with
\(1_R=1_S\).

\begin{definition}
An element \(s\in S\) is called integral over \(R\) if there exists a monic polynomial
\[
f(x)\in R[x]
\]
such that
\[
f(s)=0.
\]
\end{definition}

\begin{theorem}
For \(s\in S\), the following are equivalent:
\begin{enumerate}
\item \(s\) is integral over \(R\).
\item \(R[s]\) is contained in some finitely generated \(R\)-submodule \(M\subseteq S\)
such that
\[
sM\subseteq M.
\]
\end{enumerate}
\end{theorem}

\begin{theorem}
The elements of \(S\) which are integral over \(R\) form a subring of \(S\).
\end{theorem}

\begin{theorem}
The algebraic integers in \(\mathbb{Q}\) are precisely the ordinary integers:
\[
\{x\in \mathbb{Q}\mid x\text{ is integral over }\mathbb{Z}\}=\mathbb{Z}.
\]
\end{theorem}

\begin{theorem}
Let \(G\) be a finite group, let \(g\in G\), and let \(\chi\) be a complex character of \(G\).
Then
\[
\chi(g)
\]
is an algebraic integer.
\end{theorem}

\begin{remark}
Recall that an element of \(\C\) is called an algebraic integer if it is integral over \(\mathbb{Z}\).
\end{remark}

\section{Central Elements in Group Rings}

Consider the group rings
\[
\mathbb{C}G
\qquad\text{and}\qquad
\mathbb{Z}G.
\]
We identify
\[
\mathbb{Z}\cdot 1_G\subseteq Z(\mathbb{Z}G).
\]

\begin{proposition}
The center \(Z(\mathbb{Z}G)\) is integral over \(\mathbb{Z}\).
\end{proposition}

\begin{proof}
Let \(C_1,\dots,C_r\) be the conjugacy classes of \(G\), and define the class sums
\[
\overline C_i=\sum_{g\in C_i} g.
\]
Then
\[
Z(\mathbb{Z}G)=\bigoplus_{i=1}^r \mathbb{Z}\,\overline C_i.
\]
Thus \(Z(\mathbb{Z}G)\) is a finitely generated abelian group, hence a finitely generated
\(\mathbb{Z}\)-module. Therefore every element of \(Z(\mathbb{Z}G)\) is integral over \(\mathbb{Z}\).
\end{proof}

\begin{corollary}
Let \(C_1,\dots,C_r\) be the conjugacy classes of \(G\), and let
\[
\lambda_1,\dots,\lambda_r
\]
be algebraic integers in \(\C\). Then
\[
\sum_{i=1}^r \lambda_i\overline C_i
\]
is integral over \(\mathbb{Z}\).
\end{corollary}

Let
\[
\rho_1,\dots,\rho_r
\]
be the simple complex representations of \(G\), with degrees
\[
d_1,\dots,d_r,
\]
and characters
\[
\chi_1,\dots,\chi_r.
\]
Thus
\[
\rho_i:G\to \operatorname{GL}_{d_i}(\C).
\]
Extending linearly gives
\[
\rho_i:\mathbb{C}G\to M_{d_i}(\mathbb{C}).
\]
Under the Artin--Wedderburn decomposition
\[
\mathbb{C}G\cong \bigoplus_{i=1}^r M_{d_i}(\mathbb{C}),
\]
the map \(\rho_i\) is the projection onto the \(i\)-th factor.

\begin{proposition}
Fix \(i\), and let \(x\in Z(\mathbb{C}G)\). Then
\[
\rho_i(x)=\lambda I_{d_i}
\]
for some \(\lambda\in\mathbb{C}\). In fact, if
\[
x=\sum_{g\in G} a_g g,
\]
then
\[
\lambda
=
\frac{1}{d_i}\operatorname{tr}(\rho_i(x))
=
\frac{1}{d_i}\sum_{g\in G} a_g\chi_i(g).
\]
\end{proposition}

\begin{proof}
Since \(x\in Z(\mathbb{C}G)\), for every \(y\in \mathbb{C}G\),
\[
xy=yx.
\]
Therefore, for every \(i\),
\[
\rho_i(x)\rho_i(y)=\rho_i(y)\rho_i(x).
\]
By Schur's lemma, \(\rho_i(x)\) is a scalar matrix:
\[
\rho_i(x)=\lambda I_{d_i}.
\]
Taking traces gives
\[
\lambda
=
\frac{1}{d_i}\operatorname{tr}(\rho_i(x)).
\]
If
\[
x=\sum_{g\in G}a_g g,
\]
then
\[
\operatorname{tr}(\rho_i(x))
=
\operatorname{tr}\left(\sum_{g\in G}a_g\rho_i(g)\right)
=
\sum_{g\in G}a_g\operatorname{tr}(\rho_i(g))
=
\sum_{g\in G}a_g\chi_i(g).
\]
Hence
\[
\lambda
=
\frac{1}{d_i}\sum_{g\in G}a_g\chi_i(g).
\]
\end{proof}

\section{Degrees of Irreducible Complex Representations}

\begin{theorem}
Let \(G\) be a finite group, and let
\[
\rho_i:G\to \operatorname{GL}_{d_i}(\mathbb{C})
\]
be a simple complex representation of degree \(d_i\). Then
\[
d_i \mid |G|.
\]
\end{theorem}

\begin{proof}
Let \(\chi_i\) be the character of \(\rho_i\). Define
\[
x=\sum_{g\in G}\overline{\chi_i(g)}\,g\in \mathbb{C}G.
\]
Since
\[
\overline{\chi_i(g)}=\chi_i(g^{-1}),
\]
the coefficients are constant on conjugacy classes, so \(x\in Z(\mathbb{C}G)\).

By the previous proposition,
\[
\rho_i(x)
=
\frac{1}{d_i}
\left(\sum_{g\in G}\overline{\chi_i(g)}\chi_i(g)\right)I_{d_i}.
\]
By the orthogonality relations,
\[
\sum_{g\in G}\overline{\chi_i(g)}\chi_i(g)=|G|.
\]
Therefore
\[
\rho_i(x)=\frac{|G|}{d_i}I_{d_i}.
\]

For every \(g\in G\), the number \(\chi_i(g)\) is an algebraic integer. Hence
\[
\overline{\chi_i(g)}
=
\chi_i(g^{-1})
\]
is also an algebraic integer. Therefore
\[
x
=
\sum_{C\in \operatorname{Cl}(G)}
\overline{\chi_i(g_C)}
\left(\sum_{g\in C}g\right)
\]
is integral over \(\mathbb{Z}\), where \(g_C\) is any representative of the conjugacy class \(C\).

Thus \(\rho_i(x)\) is integral over \(\mathbb{Z}\). Hence the scalar
\[
\frac{|G|}{d_i}
\]
is integral over \(\mathbb{Z}\). But
\[
\frac{|G|}{d_i}\in \mathbb{Q}.
\]
Since the only rational algebraic integers are the ordinary integers, we get
\[
\frac{|G|}{d_i}\in \mathbb{Z}.
\]
Therefore
\[
d_i\mid |G|.
\]
\end{proof}

\section{Character Table Examples}

We recall several standard facts about irreducible characters.

\begin{itemize}
\item Linear characters of \(G\) are the irreducible characters of the abelianization
\[
G/[G,G].
\]

\item The number of irreducible characters equals the number of conjugacy classes:
\[
|\operatorname{Irr}(G)|=|\operatorname{Cl}(G)|.
\]

\item The sum of the squares of the degrees is the order of the group:
\[
|G|
=
\sum_{\chi\in \operatorname{Irr}(G)} \chi(1)^2.
\]

\item For every \(\chi\in \operatorname{Irr}(G)\),
\[
\chi(1)\mid |G|.
\]

\item The irreducible characters satisfy the usual orthogonality relations.
\end{itemize}

\begin{example}[The character table of \(S_3\)]
For \(S_3\), we have
\[
S_3/[S_3,S_3]=S_3/A_3\cong C_2.
\]
Thus \(S_3\) has two linear characters. Since
\[
|S_3|=6
\]
and the irreducible degrees satisfy
\[
6=\sum_{\chi\in\operatorname{Irr}(S_3)}\chi(1)^2,
\]
the remaining irreducible character has degree \(2\).

The conjugacy classes are
\[
(1),\qquad (12),\qquad (123).
\]
The character table is
\[
\begin{array}{c|ccc}
 & (1) & (12) & (123)\\
\hline
1 & 1 & 1 & 1\\
\varepsilon & 1 & -1 & 1\\
\chi & 2 & 0 & -1
\end{array}
\]
\end{example}

\begin{example}[The character table of \(D_8\)]
Let
\[
D_8=\langle r,s\mid r^4=s^2=1,\ srs=r^{-1}\rangle
\]
be the dihedral group of order \(8\).

One standard ordering of the conjugacy classes is
\[
C_1=\{1\},\qquad
C_2=\{s,r^2s\},\qquad
C_3=\{rs,r^3s\},
\]
\[
C_4=\{r,r^3\},\qquad
C_5=\{r^2\}.
\]
The character table is
\[
\begin{array}{c|ccccc}
 & C_1 & C_2 & C_3 & C_4 & C_5\\
\hline
\chi_1 & 1 & 1 & 1 & 1 & 1\\
\chi_2 & 1 & -1 & 1 & -1 & 1\\
\chi_3 & 1 & 1 & -1 & -1 & 1\\
\chi_4 & 1 & -1 & -1 & 1 & 1\\
\chi_5 & 2 & 0 & 0 & 0 & -2
\end{array}
\]
\end{example}
